```latex
\documentclass[11pt]{article}

\usepackage{amsmath,amssymb,amsthm}
\usepackage[margin=1in]{geometry}

\newtheorem{theorem}{Theorem}[section]
\newtheorem{lemma}[theorem]{Lemma}
\newtheorem{proposition}[theorem]{Proposition}
\newtheorem{corollary}[theorem]{Corollary}
\theoremstyle{definition}
\newtheorem{definition}[theorem]{Definition}
\theoremstyle{remark}
\newtheorem{remark}[theorem]{Remark}

\newcommand{\op}{{\mathrm{op}}}
\newcommand{\Sh}{\operatorname{Sh}}
\newcommand{\PSh}{\operatorname{PSh}}
\newcommand{\Lan}{\operatorname{Lan}}
\newcommand{\colim}{\operatorname{colim}}
\newcommand{\id}{\mathrm{id}}
\newcommand{\N}{\mathbb N}
\newcommand{\Z}{\mathbb Z}

\begin{document}

\title{A counterexample to the proposed derived unit isomorphism}
\author{}
\date{}
\maketitle

\section*{Answer}

The answer is \emph{no}.  Already for \(P=\N\), the left Kan
extension used in the definition of \(\pi^{-1}\) can discard all the
values of a presheaf except its value at the zero monoid.  We give an
explicit presheaf \(F\) on \(B_{\N}^{\op}\) for which
\[
\pi^{-1}F=0
\quad\text{but}\quad F\neq 0.
\]
Consequently
\[
F\longrightarrow \mathrm R\pi_{\ast}\pi^{-1}F
\]
cannot be a quasi-isomorphism.

The standard definition of valuative monoids used here agrees with the
usual formulation, namely integrality together with the condition that
for every \(m\in M^{\mathrm{gp}}\), either \(m\in M\) or \(-m\in M\);
compare
[pub.math.leidenuniv.nl](https://pub.math.leidenuniv.nl/~holmesdst/notes_from_talks/valuativisations_talk.pdf).

\section{The zero monoid lies in \(B_{\N}\)}

Let
\[
P=\N
\]
under addition.  It is integral, saturated, sharp, commutative, and
unital.

Write \(0\) for the one-element monoid.  It is also integral,
saturated, and sharp.

\begin{lemma}\label{lem:zero-in-B}
The object
\[
t=(\N\longrightarrow 0)
\]
belongs to \(B_{\N}\).
\end{lemma}

\begin{proof}
Take \(x=1\in\N\).  Localization at \(1\) gives
\[
\N[-1]\simeq \Z.
\]
Every element of \(\Z\) is a unit of the additive monoid \(\Z\), so
\[
\Z^\sharp=\Z/\Z^\times=0.
\]
Hence
\[
\N[-1]^\sharp\simeq 0,
\]
and the canonical localization map \(\N\to\N[-1]^\sharp\) is the
unique map \(\N\to0\).  Thus \(t\in B_{\N}\).
\end{proof}

Also, the identity object
\[
e=(\N\xrightarrow{\id}\N)
\]
belongs to \(B_{\N}\): localization at \(0\) does not change a monoid,
and \(\N^\sharp=\N\).

\section{The left Kan extension sees only \(F(t)\)}

Put
\[
\mathcal A=A_{\N},\qquad \mathcal B=B_{\N},
\]
and let
\[
i\colon \mathcal B\hookrightarrow \mathcal A
\]
be the full inclusion.  The functor in the statement is
\[
i^\op\colon\mathcal B^\op\longrightarrow\mathcal A^\op.
\]

Recall that a presheaf on \(\mathcal B^\op\) is a functor
\[
F\colon(\mathcal B^\op)^\op=\mathcal B\longrightarrow \mathbf{Ab}.
\]
For \(q\in\mathcal A^\op\), the pointwise left Kan extension is
\[
(\Lan_{i^\op}F)(q)
=
\colim_{(d,\alpha)\in(i^\op\downarrow q)}F(d),
\]
where an object of the comma category consists of \(d\in\mathcal B\)
and an arrow
\[
\alpha\colon i^\op(d)\longrightarrow q
\]
in \(\mathcal A^\op\).

\begin{lemma}\label{lem:terminal}
The object \(t=(\N\to0)\) is terminal in \(\mathcal A\), and hence
initial in \(\mathcal A^\op\).
\end{lemma}

\begin{proof}
For every monoid \(Q\), there is a unique monoid homomorphism
\(Q\to0\).  Therefore for every object \(q=(\N\to Q)\) of
\(\mathcal A\), there is a unique morphism \(q\to t\) in
\(\mathcal A\), given by \(Q\to0\).  Thus \(t\) is terminal in
\(\mathcal A\), and therefore initial in \(\mathcal A^\op\).
\end{proof}

\begin{lemma}\label{lem:initial-comma}
For every \(q\in\mathcal A^\op\), the object
\[
\bigl(t,\;i^\op(t)\to q\bigr)
\]
is initial in the comma category \((i^\op\downarrow q)\).
\end{lemma}

\begin{proof}
By Lemma~\ref{lem:terminal}, there is a unique arrow
\(i^\op(t)\to q\) in \(\mathcal A^\op\).

Let \((d,\alpha\colon i^\op(d)\to q)\) be any object of the comma
category.  Since \(t\) is terminal in \(\mathcal A\), there is a
unique morphism
\[
i(d)\longrightarrow i(t)
\]
in \(\mathcal A\), namely the map to the zero monoid.  Reversing the
arrow gives a unique morphism
\[
i^\op(t)\longrightarrow i^\op(d)
\]
in \(\mathcal A^\op\).  The triangle
\[
i^\op(t)\longrightarrow i^\op(d)\xrightarrow{\alpha}q
\]
commutes because \(t\) is initial in \(\mathcal A^\op\).  Hence the
object indexed by \(t\) maps uniquely to every object of the comma
category, so it is initial.
\end{proof}

\begin{lemma}\label{lem:initial-colimit}
Let \(J\) be a category with an initial object \(j_0\), and let
\(G\colon J\to\mathbf{Ab}\) be a functor.  Then the canonical map
\[
G(j_0)\longrightarrow\colim_JG
\]
is an isomorphism.
\end{lemma}

\begin{proof}
The maps \(G(j_0)\to G(j)\) induced by the unique morphisms
\(j_0\to j\) form a cocone.  Given any cocone
\(\{\psi_j\colon G(j)\to X\}\), compatibility forces
\[
\psi_j\circ G(j_0\to j)=\psi_{j_0}.
\]
Thus every cocone is uniquely determined by \(\psi_{j_0}\), proving
the universal property.
\end{proof}

\begin{proposition}\label{prop:Lan-constant}
For every presheaf \(F\) on \(B_{\N}^{\op}\),
\[
\Lan_{i^\op}F
\]
is the constant presheaf on \(A_{\N}^{\op}\) with value \(F(t)\).
\end{proposition}

\begin{proof}
By Lemma~\ref{lem:initial-comma}, the indexing category for
\((\Lan_{i^\op}F)(q)\) has initial object indexed by \(t\).  By
Lemma~\ref{lem:initial-colimit},
\[
(\Lan_{i^\op}F)(q)\simeq F(t).
\]
These isomorphisms are natural in \(q\), since the initial objects
and the comparison maps are unique.
\end{proof}

\section{An explicit nonzero presheaf killed by \(\pi^{-1}\)}

Define a functor
\[
F\colon B_{\N}\longrightarrow\mathbf{Ab}
\]
by
\[
F(e)=\Z,\qquad F(t)=0,
\]
and \(F(s)=0\) for every object \(s\in B_{\N}\) other than \(e\) and
\(t\).  Send every identity morphism to the corresponding identity map
and every nonidentity morphism to the zero homomorphism.

\begin{lemma}
The assignment \(F\) is a functor, hence a presheaf of abelian groups
on \(B_{\N}^{\op}\).
\end{lemma}

\begin{proof}
Identities are sent to identities by definition.  A composite
involving a nonidentity morphism is sent to a composite of zero maps,
hence to zero.  Thus composition is preserved.  Since a presheaf on
\(B_{\N}^{\op}\) is a covariant functor on \(B_{\N}\), this proves the
claim.
\end{proof}

In particular, the map associated to the unique localization morphism
\(e\to t\) is the zero map
\[
\Z=F(e)\longrightarrow F(t)=0.
\]

\begin{proposition}\label{prop:pi-inverse-zero}
For the presheaf \(F\) above,
\[
\pi^{-1}F=0.
\]
\end{proposition}

\begin{proof}
By Proposition~\ref{prop:Lan-constant},
\[
\Lan_{i^\op}F
\]
is the constant presheaf with value
\[
F(t)=0.
\]
Hence \(\Lan_{i^\op}F=0\).  By construction, \(\pi^{-1}F\) is the
birhizomic sheafification of this left Kan extension.  Sheafification
of the zero presheaf is the zero sheaf, so
\[
\pi^{-1}F=0.
\]
\end{proof}

\begin{theorem}
For \(P=\N\), there exists a presheaf \(F\) of abelian groups on
\(B_P^{\op}\) such that
\[
F\longrightarrow\mathrm R\pi_{\ast}\pi^{-1}F
\]
is not a quasi-isomorphism.
\end{theorem}

\begin{proof}
Let \(F\) be the presheaf constructed above.  By
Proposition~\ref{prop:pi-inverse-zero},
\[
\pi^{-1}F=0.
\]
Therefore
\[
\mathrm R\pi_{\ast}\pi^{-1}F
=
\mathrm R\pi_{\ast}(0)
=
0.
\]
Thus the derived unit morphism is
\[
F\longrightarrow 0.
\]
Since the topology on \(B_{\N}^{\op}\) is trivial, its category of
sheaves is its category of presheaves.  The object \(F\), regarded as
a complex concentrated in degree \(0\), satisfies
\[
H^0(F)(e)=F(e)=\Z\neq0,
\]
whereas \(H^0(0)(e)=0\).  Hence \(F\to0\) is not a
quasi-isomorphism.
\end{proof}

\begin{corollary}
The asserted statement is false in general.
\end{corollary}

\begin{remark}
The failure occurs before sheafification and before any derived
functor is applied: the terminal localization \(\N\to0\) lies in
\(B_{\N}\), and the pointwise left Kan extension along
\(B_{\N}^{\op}\hookrightarrow A_{\N}^{\op}\) is constant with value at
that terminal localization.  Consequently, no property of the
birhizomic topology can repair the loss of the value \(F(\N)\) for the
particular presheaf above.
\end{remark}

\end{document}
```